\newpage
\chapter{活動環境と担当の引継ぎ}

\section{場や権限を変更するとき}

新しい場を設ける担当者は、用途、対象者、管理する者を決めて案内する。
既存の場を整理するときは、学生の利用状況と残すべき活動の記録を確認し、影響を受ける者に変更の内容と時点を知らせる。
閲覧・編集権限は業務に必要な範囲で付与し、役職の技術的な権限を学生間の序列として扱わない。
場の案内だけで規則にない禁止や処分事由を作らない。

\section{設定と自動化の点検}

設定を変更した担当者は、目的、変更内容、確認した動作及び戻し方を記録する。
反復する事務は自動化できるが、学生の資格、認可、処分その他の権利利益に影響する結論を自動処理の結果だけで確定しない。
意図しない結果が生じたときの連絡先と訂正方法を、担当の教務主事が把握する。

\section{担当交代}

交代する者は、継続中の申請・問合せ、判断の理由、公示済みの案内、場の設定及び記録の所在を後任へ伝える。
後任は、引き継いだ権限で必要な業務ができるか確認し、不要になった前任者の権限を整理する。
履修登録のフォームと回答シートは第４章の手順で引き継ぐ。
引継ぎで判明した未決定事項を、前任者の慣行だけで制度上の決定とみなさない。

\bigskip
\noindent

この手引きは、２０２６年９月１日から施行する。

\bigskip
\noindent 沿革
\begin{itemize}
\item ２０２６年８月２９日制定
\item ２０２６年８月３０日改正
\item ２０２６年１０月１日改正
\end{itemize}
